\documentclass[journal,twoside,web]{ieeecolor}
\usepackage{generic}
\usepackage{cite}
\usepackage{amsmath,amssymb,amsfonts}
\usepackage{algorithmic}
\usepackage{graphicx}
\usepackage{textcomp}
\usepackage{longtable}   % for long tables across pages
\usepackage{multirow}    % for multirow cells
\usepackage{array}       % better column formatting
\usepackage{booktabs}    % optional (only if you use \toprule etc.)
\def\BibTeX{{\rm B\kern-.05em{\sc i\kern-.025em b}\kern-.08em
    T\kern-.1667em\lower.7ex\hbox{E}\kern-.125emX}}
\markboth{\journalname, VOL. XX, NO. XX, XXXX}
{Author \MakeLowercase{\textit{et al.}}: Title}
\begin{document}
\title{Preparation of Papers for IEEE Trans on Industrial Informatics (February 2022)}
\author{First A. Author, \IEEEmembership{Fellow, IEEE}, Second B. Author, and Third C. Author, Jr., \IEEEmembership{Member, IEEE}
\thanks{This work was supported by the \textit{SecureIoTM: Ultra-low-energy IoT Intrusion Detection Systems using Logic-based Tsetlin Machines} project (Grant Number: 342167), funded by the Research Council of Norway under the TEKNOLOGIKONVER program. 
The authors acknowledge the support from the Center of Artificial Intelligence Research at the University of Agder, Sørlandet Hospital Trust, and Newcastle University for their collaboration in this interdisciplinary research initiative.}% <-this % stops a space
\thanks{Rakesh Reddy Yakkati, Per-Arne Andersen, Lei Jiao, and Ole-Christoffer Granmo are with the Department of ICT, 
University of Agder, Grimstad 4879, Norway (\{rakeshry; per.andersen; lei.jiao; ole.granmo\}@uia.no). }
\thanks{Rishad Shafik is with the Department of Microelectronic Systems, Newcastle University, Newcastle upon Tyne, UK (e-mail: rishad.shafik@newcastle.ac.uk). (Corresponding authors: Per-Arne Andersen and Lei Jiao).}}

\maketitle

\begin{abstract}
Intrusion Detection Systems (IDS) play a critical role in securing heterogeneous Internet of Things (IoT) infrastructures, including Industrial IoT (IIoT), Internet of Medical Things (IoMT), and cyber-physical systems. However, stringent constraints on computation, memory, and latency hinder the deployment of conventional machine learning models in edge computing environments. This paper presents a lightweight and deployment-efficient IDS framework designed for resource-constrained edge systems, emphasizing efficient booleanization, a compact Tsetlin Machine architecture, and low-latency inference.
An efficient binarization strategy termed GLADE is proposed to transform heterogeneous data streams into compact and discriminative Boolean representations while reducing redundancy and controlling feature expansion. In contrast to thermometric encoding, GLADE generates compact feature spaces that improve learning efficiency and reduce memory footprint. The GLADE booleanizer is integrated with a Fuzzy Pattern Tsetlin Machine (FPTM), enabling robust and interpretable rule-based learning under uncertainty while preserving a compact model structure. This integration reduces clause complexity and enhances generalization across diverse intrusion detection scenarios.
The proposed framework is evaluated on benchmark datasets including WUSTL-EHMS-2020, NSL-KDD, TON\_IoT, and MedSec-25. Experimental results demonstrate competitive detection performance compared to conventional machine learning models. Model storage remains compact with models serialized in two formats: a non-compressed pickle format for Python-based deployment and a compressed Zstd (level-1) format that supports efficient storage and cross-platform interoperability. On a resource-constrained edge platform such as the Raspberry Pi 5, GLADE binarisation runs up to 19× faster than the standard Tsetlin Machine binariser and up to 11× faster than thermometric encoding without compromising detection accuracy.
\end{abstract}

% FIX 1: IEEEkeywords entries must be comma-separated
\begin{IEEEkeywords}
Intrusion Detection Systems,
Internet of Things,
Edge Computing,
Tsetlin Machine,
Booleanization,
Lightweight Machine Learning,
Cyber-Physical Systems
\end{IEEEkeywords}

\section{Introduction}
\label{sec:introduction}

\IEEEPARstart{T}{he} Internet of Things (IoT) has evolved into a large-scale cyber-physical infrastructure comprising billions of connected sensors, controllers, and monitoring systems deployed across industrial, healthcare, and critical environments \cite{hussain2020ml, algaradi2020survey}. These systems enable tight integration between digital intelligence and physical processes, but also significantly increase exposure to cyber threats. A compromised node in an industrial control system can disrupt production and incur substantial economic loss, while attacks in healthcare environments can directly impact patient safety and system reliability \cite{zhukabayeva2025cybersecurity}. Consequently, intrusion detection has become a fundamental requirement for secure IoT operation.

Machine learning has been widely adopted for intrusion detection, leading to substantial improvements in detection accuracy on benchmark datasets \cite{kikissagbe2024review, jouhari2024lightweight}. However, a persistent challenge remains. The environments where intrusion detection is most critical are also those with the most severe resource constraints. Edge devices such as industrial controllers and embedded medical systems operate under strict limits on memory, computation, and latency \cite{wang2023lightweight}. Models that perform well under server-side conditions are often impractical when deployed on such hardware \cite{kundu2025review}. This mismatch between detection capability and deployment feasibility remains a central challenge in IoT security.

Efforts to address this issue have primarily focused on reducing the cost of existing models through compression techniques such as
quantization, pruning, and knowledge distillation \cite{liu2025compression}. An alternative approach is to reconsider the feature representation itself. By transforming input features into Boolean form, inference can be performed using bit-level operations,
eliminating the need for floating-point computation \cite{bacellar2017thermometer}. However, conventional encoding methods,
such as thermometric encoding, lead to significant feature expansion, which increases model size and learning complexity \cite{nguyen2025binarized}.

Tsetlin Machines (TMs) provide a rule-based learning framework that operates directly on Boolean inputs using propositional logic clauses
\cite{granmo2018tsetlin}. Inference relies solely on logical operations, making TMs well-suited for resource-constrained hardware.
The Fuzzy Pattern Tsetlin Machine (FPTM) further improves robustness by allowing partial clause matching, which is beneficial in noisy and
heterogeneous IoT environments \cite{hnilov2025fptm}. Despite these advantages, the efficiency of TM-based systems is strongly influenced
by the quality of the input encoding. Suboptimal booleanization leads to unnecessarily large clause sets and limits achievable efficiency.

The principal contributions of this work are as follows:

\begin{itemize}
    \item \textbf{GLADE Booleanizer:} A novel feature booleanization method that generates compact and discriminative Boolean representations, reducing feature expansion and clause complexity compared to thermometric encoding.

    \item \textbf{GLADE-FPTM Integration:} A unified framework that combines GLADE with the FPTM to achieve robust, interpretable, and
    memory-efficient intrusion detection.

    \item \textbf{Comprehensive Evaluation:} Experimental validation across multiple heterogeneous IoT intrusion detection datasets,
    demonstrating competitive detection accuracy with reduced model size and computational cost.

    \item \textbf{Edge Deployment Analysis:} Empirical evaluation on a resource-constrained platform, demonstrating significant
    improvements in inference efficiency and real-time feasibility.
\end{itemize}

The remainder of this paper is organized as follows. Related work is reviewed in Section~\ref{sec:related}. The proposed GLADE method is
described in Section~\ref{sec:glade}. The FPTM framework is presented in Section~\ref{sec:fptm}. Experimental setup and datasets are
introduced in Section~\ref{sec:experiments}. Results are discussed in Section~\ref{sec:results}, and conclusions are drawn in Section~\ref{sec:conclusion}.


\section{Related Work}
\label{sec:related}

% ---------------------------------------------------------------
\subsection{Machine Learning for IoT and IIoT Intrusion Detection}
% ---------------------------------------------------------------

Machine learning has been widely applied to intrusion detection in
IoT environments, encompassing signature-based, anomaly-based, and
hybrid approaches. A recurring challenge identified in prior work is
the difficulty of deploying such models on resource-constrained edge
devices \cite{hussain2020ml, algaradi2020survey}. Deep learning
methods have demonstrated strong performance in industrial and
medical IoT settings, particularly on benchmark datasets
\cite{wang2023lightweight, jouhari2024lightweight, hady2020iomt}.

Benchmark datasets such as NSL-KDD and TON\_IoT have enabled
systematic evaluation of intrusion detection methods
\cite{tavallaee2009nslkdd, moustafa2021toniot}. However, most studies
focus on server-side evaluation, with limited consideration of
deployment constraints such as memory usage and inference latency.

% ---------------------------------------------------------------
\subsection{Lightweight and Edge-Deployable IDS}
% ---------------------------------------------------------------

To address deployment challenges, lightweight IDS solutions have been
developed using model compression techniques, including pruning,
quantization, and knowledge distillation \cite{liu2025compression}.
Hybrid architectures combining convolutional and recurrent components
have also been proposed for embedded platforms
\cite{jouhari2024lightweight}. While these approaches reduce resource
requirements, they are typically applied as post-hoc optimizations to
models designed for floating-point computation. Consequently, they
introduce approximation error and do not eliminate reliance on
floating-point arithmetic.

% ---------------------------------------------------------------
\subsection{Tsetlin Machines, Booleanization, and Edge Security}
% ---------------------------------------------------------------

The Tsetlin Machine (TM) is a rule-based learning framework that
performs classification using propositional logic clauses constructed
by teams of Tsetlin Automata \cite{granmo2018tsetlin}. Its reliance on
bitwise operations makes it inherently suitable for deployment on
resource-constrained hardware \cite{kundu2025review}. Competitive
performance has been demonstrated across multiple intrusion detection
benchmarks while maintaining interpretability and low inference cost
\cite{abeyrathna2020intrusion, gunvaldsen2023iot, jaiswal2025iomt}.

The Fuzzy Pattern Tsetlin Machine (FPTM) extends this framework by
allowing partial clause matching, improving robustness to noisy and
heterogeneous data \cite{hnilov2025fptm}.

A key requirement for TM-based learning is the transformation of
features into Boolean form. Thermometric encoding is widely used for
this purpose, but its dimensionality grows linearly with the number of
thresholds, leading to significant feature expansion and increased
model complexity \cite{granmo2018tsetlin, gunvaldsen2023iot}.
Alternative binarization methods proposed for neural networks rely on
gradient-based optimization and are not directly applicable to TM
learning \cite{nguyen2025binarized}.

Despite these advances, the booleanization process has remained
largely decoupled from the learning algorithm. Its impact on clause
complexity, memory usage, and inference efficiency has not been
systematically addressed.


%% ========================================================================
% FIX 3: GLADE section moved before FPTM to match the intro roadmap
\section{GLADE Binariser}
\label{sec:glade}

Tsetlin Machines require Boolean inputs; therefore, transforming raw
numerical features into informative binary representations is a
critical preprocessing step. Conventional discretisation methods,
such as equal-width and equal-frequency binning, apply a fixed
allocation of thresholds across all features, which often leads to
inefficient use of representational capacity and poorly placed
thresholds in low-density regions. To address these limitations, the
Gap-aware Lightweight Adaptive Discretisation Engine (GLADE) is
proposed as a core contribution of this work.

GLADE generates Boolean literals through threshold comparisons of
the form
\begin{equation}
  b_k = \mathbf{1}\!\left[ x_{f_k} \geq \tau_k \right],
  \quad k = 1, \ldots, K,
\end{equation}
where $x_{f_k}$ denotes the value of feature $f_k$, $\tau_k$ is a
learned threshold, and $K$ is the total number of generated bits.
The method is governed by a single hyperparameter,
$n_\text{bins}$, which defines the maximum number of thresholds per
feature. The binarisation process consists of four stages.

%% ------------------------------------------------------------------------
\subsection{Step 1: Adaptive Bit Allocation}
\label{subsec:budget}

An individual bit budget $b_j$ is assigned to each feature based on
its statistical properties.

\textbf{Categorical features:} When the number of distinct values
does not exceed $n_\text{bins}$, the feature is treated as
categorical, and thresholds are placed between consecutive values,
yielding $|\mathcal{U}_j| - 1$ bits.

\textbf{Sparse features:} For features with a high proportion of
zero values, allocating the full budget is inefficient. When the
zero fraction $\rho_j$ exceeds 0.3, the effective number of
informative values is estimated as
\begin{equation}
  d_j = \max(1 - \rho_j,\ 0.01),
\end{equation}
\begin{equation}
  \hat{e}_j = \max\!\left(\tilde{n}_j \cdot d_j^{2},\ 2\right),
\end{equation}
\begin{equation}
  b_j = \max\!\left(1,\ \min\!\left(
        \lceil \log_2 \hat{e}_j \rceil + 2,\;
        n_\text{bins}\right)\right),
\end{equation}
where $\tilde{n}_j$ denotes the number of distinct non-zero values.
For dense features, the full budget $n_\text{bins}$ is retained.

%% ------------------------------------------------------------------------
\subsection{Step 2: Gap-Aware Threshold Placement}
\label{subsec:gaps}

Thresholds are initially estimated using equal-frequency quantiles.
However, this strategy becomes suboptimal when the feature
distribution contains large empty intervals. To address this issue,
GLADE detects structural gaps by comparing the maximum spacing
$g_\text{max}$ between consecutive sorted values with the median
spacing $\tilde{g}$. When
\begin{equation}
  g_\text{max} > 5\,\tilde{g},
\end{equation}
the distribution is considered discontinuous, and candidate
thresholds are relocated to the midpoints of the detected gaps.
This ensures that thresholds separate regions with actual data
support.

%% ------------------------------------------------------------------------
\subsection{Step 3: Entropy-Based Pruning}
\label{subsec:entropy}

Each threshold is evaluated using the entropy of the resulting
binary split,
\begin{equation}
  H(\tau) = -p_\tau \log p_\tau - (1-p_\tau)\log(1-p_\tau),
\end{equation}
where $p_\tau = \Pr[x_j \geq \tau]$. Thresholds with
$H(\tau) < 0.05$ nats are considered uninformative and are removed.
To maintain the allocated budget, new thresholds are introduced by
iteratively splitting the widest remaining intervals. This procedure
ensures that the available capacity is concentrated in regions where
the feature exhibits meaningful variation.

%% ------------------------------------------------------------------------
\subsection{Step 4: Local Refinement}
\label{subsec:refine}

A local refinement step is applied to each remaining threshold.
Three candidate positions are evaluated,
\begin{equation}
  \tau^{-} = \tau - \tfrac{1}{4}(\tau - \tau_\text{left}), \quad
  \tau^{0} = \tau, \quad
  \tau^{+} = \tau + \tfrac{1}{4}(\tau_\text{right} - \tau),
\end{equation}
where $\tau_\text{left}$ and $\tau_\text{right}$ denote adjacent
thresholds. The position that maximises the Bernoulli variance
$p(1-p)$ is selected, promoting balanced and informative splits.

%% ------------------------------------------------------------------------
\subsection{Inference}
\label{subsec:inference}

During inference, binarisation is performed as a set of independent
comparisons between feature values and stored thresholds. The
resulting Boolean outputs can be packed into compact bit vectors,
enabling efficient memory usage. The computational complexity scales
linearly with $K$, and the absence of branching or dynamic memory
allocation makes the method suitable for deployment on embedded and
resource-constrained platforms.


%% ========================================================================
% FIX 4: FPTM section moved after GLADE to match the intro roadmap
\section{The Fuzzy-Pattern Tsetlin Machine}
\label{sec:fptm}
This paper employs the Fuzzy-Pattern Tsetlin Machine
(FPTM)~\cite{hnilov2025fuzzy} as its classifier. The public
reference implementation is used without modification to any part
of the training procedure. The following subsections describe the
forward-pass computation in sufficient detail for the analysis in
Section~\ref{sec:results} to be self-contained.
% FIX 5: \ref{sec:hardware} replaced with \ref{sec:results}
%        (no hardware-analysis section exists in this draft)

\subsection{Clause Representation and Mismatch Counting}
\label{sec:fptm-clause}
For each output class $k \in \{1,\ldots,K\}$ and each polarity
$\sigma \in \{+,-\}$, the FPTM stores a bank of $C$ \emph{clauses}.
Every clause encodes a conjunctive pattern over a binarised input
vector $\mathbf{x} \in \{0,1\}^{N}$.  During training, Tsetlin
automata decide, for each feature and its complement, whether the
corresponding literal is \emph{active} in the clause.  At inference
time this decision is frozen into two packed bitmasks of
$H = \lceil N/64 \rceil$ 64-bit words:
$\mathbf{w}^{I}_{\phi}$, whose set bits identify the positive
literals active in clause $\phi$, and
$\mathbf{w}^{E}_{\phi}$, whose set bits identify the active
negated literals.

Unlike the classical Tsetlin Machine, which disqualifies a clause
the moment any single active literal disagrees with the input, the
FPTM quantifies the degree of disagreement.  The
\emph{mismatch count} for clause $\phi$ on input $\mathbf{x}$ is:
%
\begin{equation}
  m(\phi,\mathbf{x})
  \;=\;
  \operatorname{popcount}\!\bigl(
      \mathbf{w}^{I}_{\phi} \wedge \neg\mathbf{w}_{\mathbf{x}}
  \bigr)
  \;+\;
  \operatorname{popcount}\!\bigl(
      \mathbf{w}^{E}_{\phi} \wedge \mathbf{w}_{\mathbf{x}}
  \bigr),
  \label{eq:mismatch}
\end{equation}
%
where $\mathbf{w}_{\mathbf{x}}$ is the packed representation of
$\mathbf{x}$. The first popcount detects active positive literals
that are absent from the input; the second detects active negated
literals that are present. Their sum therefore equals the total
number of active literals whose value conflicts with $\mathbf{x}$.

\subsection{Graded Output and Per-Clause Capacity}
\label{sec:fptm-graded}
Each clause is assigned a capacity
%
\begin{equation}
  \kappa_{\phi} \;=\; \min\!\bigl(|\phi|,\;\mathrm{LF}\bigr),
  \label{eq:kappa}
\end{equation}
%
where $|\phi|$ is the number of active literals in clause $\phi$ and
$\mathrm{LF}$ is the \emph{Literal Failures} hyperparameter
of~\cite{hnilov2025fuzzy}, fixed at inference. The graded output of
the clause is then:
%
\begin{equation}
  o(\phi,\mathbf{x})
  \;=\; \max\!\bigl(\kappa_{\phi} - m(\phi,\mathbf{x}),\;0\bigr).
  \label{eq:graded}
\end{equation}
%
When no active literal conflicts with the input,
$o(\phi,\mathbf{x}) = \kappa_{\phi}$, the maximum possible
contribution. Each additional conflict reduces the output by one,
and once the number of conflicts reaches $\kappa_{\phi}$ the clause
contributes nothing. Because $\kappa_{\phi}$ is bounded by
$|\phi|$, a clause with fewer active literals than $\mathrm{LF}$
cannot claim votes it has not earned.

This linear decay has two practical consequences. First, a single
clause can respond meaningfully to a neighbourhood of inputs rather
than a single binary point, which substantially reduces the total
number of clauses a model requires. Second, isolated feature
corruptions lower the clause output only slightly, giving the model
inherent tolerance to noisy inputs.

\subsection{Class Vote and Prediction}
\label{sec:fptm-vote}
Two disjoint sets of clauses are maintained per class: $\Phi_{k}^{+}$
accumulates evidence \emph{for} class $k$, while $\Phi_{k}^{-}$
accumulates evidence \emph{against} it. The signed vote total is:
%
\begin{equation}
  V_{k}(\mathbf{x})
  \;=\;
  \sum_{\phi \in \Phi_{k}^{+}} o(\phi,\mathbf{x})
  \;-\;
  \sum_{\phi \in \Phi_{k}^{-}} o(\phi,\mathbf{x}),
  \label{eq:vote}
\end{equation}
%
and the predicted label is $\hat{k} = \arg\max_{k}\,V_{k}(\mathbf{x})$.
All $2CK$ clauses read the same read-only input bitmasks and produce
independent scalar outputs, so the computation partitions without
inter-clause communication and maps directly onto data-parallel
hardware.
\label{sec:tm-architecture}
\begin{figure*}[h]
    \centering    \includegraphics[width=2\columnwidth]{TII-Articles-LaTeX-template/Fptm1.pdf}
    \caption{Tsetlin Machine architecture showing the relationship between input features, Tsetlin automata, logical clauses, and classification.}
    \label{fig:tm-architecture}
\end{figure*}

\section{Experimental Setup}
\label{sec:experiments}   % FIX 6: was \label{sec:setup}; intro references sec:experiments

\subsection{Datasets}
\label{sec:datasets}

To comprehensively evaluate the proposed method, we employ four publicly
available benchmarks that cover both conventional IoT network traffic
and the more complex Internet of Medical Things (IoMT) domain. These
datasets provide diversity in scale, feature representation, and attack
taxonomy. Table~\ref{tab:datasets_transposed} summarises their key
characteristics.
% FIX 7: was \ref{tab:datasets}; correct label is tab:datasets_transposed
Record and feature counts correspond to the original datasets prior to
preprocessing, whereas training and testing splits follow the unified
pipeline described in Section~\ref{sec:preprocessing}.

\textbf{NSL-KDD}~\cite{tavallaee2009nslkdd} is a widely used benchmark
for network intrusion detection. It was introduced to address the
limitations of the KDD CUP 99 dataset, particularly redundancy and
class imbalance. The dataset contains 41 handcrafted features derived
from TCP/IP connections and categorises traffic into five classes:
Normal, DoS, Probe, R2L, and U2R. Its more challenging test set makes it
suitable for robust performance evaluation.

\textbf{TON\_IoT}~\cite{moustafa2021ton} represents realistic IoT/IIoT
environments with heterogeneous traffic collected from a distributed
testbed comprising cloud, edge, and IoT layers. It includes ten attack
categories such as DoS, DDoS, ransomware, and MITM. The dataset provides
44 features capturing protocol, statistical, and payload-level
characteristics.

\textbf{MedSec-25}~\cite{medsec25} is a recent IoMT dataset derived from
a simulated healthcare environment with interconnected medical devices.
It consists of 78 features and five classes (Normal, DoS, DDoS, MITM,
and Reconnaissance), with a total of 554,534 samples. Its rich feature
space and large scale make it representative of modern IoMT scenarios.

\textbf{WUSTL-EHMS-2020}~\cite{hady2020wustl} combines network traffic
features with physiological signals collected from wearable medical
devices in a real healthcare monitoring system. It includes three
classes (Normal, Spoofing, and Injection) and 44 features. This hybrid
dataset enables evaluation of models that integrate cyber and physical
context for intrusion detection.

\begin{table}[t]
\caption{Summary of IoT and IoMT intrusion detection datasets used in this study (transposed).}
\label{tab:datasets_transposed}
\centering
\renewcommand{\arraystretch}{1.2}
\small
\begin{tabular}{|l|c|c|c|c|}
\hline
\textbf{Attribute} & \textbf{NSL-KDD} & \textbf{TON\_IoT} & \textbf{MedSec-25} & \textbf{WUSTL-EHMS-2020} \\
\hline
Records   & 148,517 & 211,043 & 554,534 & 16,318 \\
\hline
Features  & 41      & 44      & 78      & 44 \\
\hline
Classes   & 5       & 10      & 5       & 3 \\
\hline
Domain    & IoT network flow & IoT multi-protocol & IoMT network flow & IoMT (network + biomedical) \\
\hline
Train     & 118,813 & 168,834 & 443,627 & 13,054 \\
\hline
Test      & 29,704  & 42,209  & 110,907 & 3,264 \\
\hline
\end{tabular}
\end{table}

% FIX 8: Added preprocessing subsection so \ref{sec:preprocessing} resolves
\subsection{Data Preprocessing}
\label{sec:preprocessing}

A unified preprocessing pipeline is applied across all datasets and models to ensure consistency and fairness. Raw features are first loaded and class labels are encoded into numeric form. Non-informative attributes such as identifiers are removed, and categorical features are transformed using one-hot encoding where applicable. The data is then split into training and testing sets using a stratified 80/20 partition to preserve class distributions. Data quality issues are addressed by replacing infinite values with missing values and imputing missing entries using the median computed from the training set only. To reduce redundancy and improve robustness, constant features are eliminated based on zero variance in the training data, followed by the removal of highly correlated features using a threshold of $|r| > 0.99$, also determined exclusively on the training set.

\subsection{FPTM Configuration}
\label{sec:fptm-config}

The Fuzzy Pattern Tsetlin Machine (FPTM) is configured using
dataset-specific hyperparameters obtained through tuning.
The model is parameterised by $C$ (number of clauses per class per
polarity), $T$ (voting target), $s$ (specificity), $L$ (maximum number
of literals per clause), $L_f$ (fuzzy clamp threshold), and $E$
(training epochs).

All experiments employ the proposed GLADE binariser with a fixed
discretisation parameter of $n_{\text{bins}} = 15$ across all datasets.

The hyperparameters are selected via a grid search over the
configuration space for each dataset, subject to the constraint that
$C \leq 100$ and $C$ is an even number. The final configurations are
reported in Table~\ref{tab:tm_hparams}.

\begin{table}[t]
\centering
\caption{Tuned FPTM hyperparameters per dataset.}
\label{tab:tm_hparams}
\renewcommand{\arraystretch}{1.2}
\small
\begin{tabular}{|l|c|c|c|c|c|c|}
\hline
\textbf{Dataset} & $C$ & $T$ & $s$ & $L$ & $L_f$ & $E$ \\
\hline
NSL-KDD         & 90  & 12 & 200 & 40 & 8  & 80  \\
\hline
TON\_IoT        & 100 & 15 & 20  & 25 & 8  & 200 \\
\hline
MedSec-25       & 80  & 12 & 75  & 30 & 8  & 300 \\
\hline
WUSTL-EHMS-2020 & 60  & 8  & 300 & 50 & 15 & 200 \\
\hline
\end{tabular}
\end{table}

\subsection{Baselines}
\label{sec:baselines}

Eight baseline models are evaluated under the same preprocessing
pipeline. The baseline set comprises standard machine learning
classifiers with library-default configurations, together with
variants whose capacity is aligned with the tuned FPTM.

The default baselines include XGBoost, Random Forest, Logistic
Regression, Linear SVM, Gaussian Naive Bayes, and k-Nearest Neighbours
($k=5$).

To ensure a fair comparison, additional variants are constructed by
matching model capacity to the FPTM clause count $C$. In particular,
the number of estimators is set proportional to $C$.
These variants are denoted with the suffix \texttt{\_Cmatched}. All
models are implemented using scikit-learn~\cite{pedregosa2011scikit}
and the official XGBoost Python package.
\begin{table}[t]
\centering
\caption{Baseline model configurations. Starred ($^\ast$) models have
capacity matched to the FPTM clause count $C$.}
\label{tab:model_configs}
\renewcommand{\arraystretch}{1.2}
\begin{tabular}{|l|l|}
\hline
\textbf{Model} & \textbf{Configuration} \\
\hline
\multicolumn{2}{|l|}{\textit{Default baselines}} \\
\hline
XGBoost            & n=100, max depth=6, learning rate=0.1 \\
\hline
Random Forest      & n=100, max depth=none \\
\hline
k-NN ($k=5$)       & uniform weights, Minkowski distance \\
\hline
Logistic Regression & L2 regularisation, $C=1.0$, LBFGS \\
\hline
Linear SVM         & $C=1.0$, squared hinge loss \\
\hline
Gaussian NB        & default priors \\
\hline
\multicolumn{2}{|l|}{\textit{FPTM-matched variants}} \\
\hline
XGBoost\_Cmatched$^\ast$      & n\_estimators = $C$ \\
\hline
RandomForest\_Cmatched$^\ast$ & n\_estimators = $C$ \\
\hline
\end{tabular}
\end{table}

\section{Results}
\label{sec:results}

\subsection{Binariser Comparison}
\label{sec:res-binarisers}

Table~\ref{tab:binarizer_comparison} evaluates the impact of the
binarisation scheme while keeping the FPTM hyperparameters fixed.
Three approaches are compared: GLADE, the \texttt{StandardBinarizer}
from TMU, and the equal-frequency \texttt{KBinsDiscretizer} from
scikit-learn.

Across all datasets, GLADE produces a substantially lower number of
binary features compared to \texttt{StandardBinarizer}, typically
reducing the representation size by approximately 45\%--65\%.
Despite this reduction, comparable or improved classification
performance is observed in most cases. The performance difference is
particularly pronounced on TON\_IoT and MedSec-25, where alternative
methods exhibit reduced effectiveness on sparse feature distributions.
As summarised in the final column of
Table~\ref{tab:binarizer_comparison}, GLADE delivers the best overall
trade-off---the fewest bits at the highest accuracy---on every
dataset.

\begin{table*}[t]
\centering
\renewcommand{\arraystretch}{1.2}
\small
\caption{Comparison of binarisation methods with $n_\text{bins}=15$.
Best values per dataset and metric are shown in \textbf{bold}; the
final column reports the method with the best overall trade-off
(fewest bits at the highest accuracy and macro F1).}
\label{tab:binarizer_comparison}
\begin{tabular}{|l|ccc|ccc|ccc|c|}
\hline
 & \multicolumn{3}{c|}{\textbf{GLADE}}
 & \multicolumn{3}{c|}{\textbf{StandardBinarizer}}
 & \multicolumn{3}{c|}{\textbf{KBinsDiscretizer}} & \\
\hline
\textbf{Dataset} & Bits & Acc & F1 & Bits & Acc & F1 & Bits & Acc & F1 & \textbf{Best} \\
\hline
NSL-KDD         & \textbf{253} & \textbf{0.9955} & 0.9311          & 472 & 0.9941 & \textbf{0.9342} & 448 & 0.9854 & 0.8690 & GLADE \\
\hline
TON\_IoT        & \textbf{122} & \textbf{0.9688} & \textbf{0.9385} & 186 & 0.8416 & 0.7947          & 177 & 0.6411 & 0.5663 & GLADE \\
\hline
MedSec-25       & \textbf{351} & \textbf{0.9703} & \textbf{0.8681} & 698 & 0.8830 & 0.7265          & 653 & 0.8461 & 0.5532 & GLADE \\
\hline
WUSTL-EHMS-2020 & \textbf{247} & \textbf{0.9393} & \textbf{0.7630} & 313 & 0.9381 & 0.7388          & 296 & 0.9130 & 0.6259 & GLADE \\
\hline
\end{tabular}
\end{table*}

\subsection{Main Classification Benchmark}
\label{sec:res-main}

% FIX 9: was \ref{tab:complete_results} (undefined); corrected to the
%         two actual result tables: tab:results_part1 and tab:results_part2
Tables~\ref{tab:results_part1} and~\ref{tab:results_part2} report the
classification performance of the tuned GLADE+FPTM pipeline alongside
the baseline models across the four datasets. Results are presented in
terms of accuracy, macro F1, precision, recall, serialized model size,
and mean per-sample inference time.

Across datasets, the GLADE+FPTM pipeline achieves competitive
performance relative to the baseline methods. Its macro F1 is
consistently among the top-performing models, remaining close to
strong tree-based baselines while requiring a model footprint that is
one to three orders of magnitude smaller. On WUSTL-EHMS-2020, the
pipeline ranks second in terms of macro F1, indicating its ability
to maintain balanced performance on smaller and more challenging
datasets. In terms of latency, GLADE+FPTM inference costs only a few
microseconds per sample---comparable to the fastest linear baselines
and roughly three orders of magnitude faster than $k$-NN, which must
scan the full training set at prediction time.

\begin{table*}[t]
\centering
\caption{Classification results on NSL-KDD and TON\_IoT datasets.
Inference time is the mean per-sample latency.}
\label{tab:results_part1}
\begin{tabular}{llcccccc}
\hline
\textbf{Dataset} & \textbf{Model} & \textbf{Accuracy} & \textbf{Macro F1} & \textbf{Precision} & \textbf{Recall} & \textbf{Size (KB)} & \textbf{Inference ($\mu$s/sample)} \\
\hline

\multirow{9}{*}{NSL-KDD}
 & XGBoost                & 0.9956 & 0.9488 & 0.9442 & 0.9558 & 1266 & 9.988 \\
 & XGBoost\_Cmatched      & 0.9952 & 0.9429 & 0.9419 & 0.9458 & 1161 & 9.122 \\
 & GLADE+FPTM             & 0.9951 & 0.9301 & 0.9399 & 0.9210 & 51.1 & 3.387 \\
 & kNN\_5                 & 0.9894 & 0.8855 & 0.9052 & 0.8691 & 110463 & 4118.039 \\
 & LogisticRegression     & 0.9769 & 0.8580 & 0.8877 & 0.8337 & 8.6  & 1.420 \\
 & RandomForest           & 0.9926 & 0.8563 & 0.9647 & 0.8167 & 27239 & 12.697 \\
 & LinearSVM              & 0.9655 & 0.8530 & 0.8986 & 0.8194 & 8.6  & 1.393 \\
 & RandomForest\_Cmatched & 0.9625 & 0.5843 & 0.5821 & 0.5869 & 962  & 6.124 \\
 & NaiveBayes             & 0.7567 & 0.4423 & 0.5765 & 0.6726 & 13.2 & 4.652 \\
\hline

\multirow{9}{*}{TON\_IoT}
 & XGBoost                & 0.9775 & 0.9511 & 0.9439 & 0.9614 & 2508 & 22.125 \\
 & XGBoost\_Cmatched      & 0.9764 & 0.9498 & 0.9424 & 0.9608 & 1899 & 18.138 \\
 & RandomForest           & 0.9765 & 0.9467 & 0.9366 & 0.9637 & 62383 & 11.077 \\
 & kNN\_5                 & 0.9742 & 0.9441 & 0.9430 & 0.9461 & 51444 & 5225.119 \\
 & GLADE+FPTM             & 0.9705 & 0.9418 & 0.9334 & 0.9558 & 76.2 & 8.567 \\
 & RandomForest\_Cmatched & 0.8771 & 0.7814 & 0.7961 & 0.7798 & 841  & 11.299 \\
 & LogisticRegression     & 0.7076 & 0.6347 & 0.7154 & 0.6494 & 5.5  & 0.632 \\
 & LinearSVM              & 0.7114 & 0.6273 & 0.7106 & 0.6486 & 5.5  & 0.629 \\
 & NaiveBayes             & 0.3875 & 0.3269 & 0.3694 & 0.4201 & 8.5  & 3.161 \\
\hline
\end{tabular}
\end{table*}


\begin{table*}[t]
\centering
\caption{Classification results on MedSec-25 and WUSTL-EHMS-2020 datasets.
Inference time is the mean per-sample latency; ``--'' denotes a
configuration not benchmarked.}
\label{tab:results_part2}
\begin{tabular}{llcccccc}
\hline
\textbf{Dataset} & \textbf{Model} & \textbf{Accuracy} & \textbf{Macro F1} & \textbf{Precision} & \textbf{Recall} & \textbf{Size (KB)} & \textbf{Inference ($\mu$s/sample)} \\
\hline

\multirow{9}{*}{MedSec-25}
 & XGBoost                & 0.9768 & 0.8984 & 0.9005 & 0.9004 & 1701 & 12.091 \\
 & RandomForest           & 0.9767 & 0.8980 & 0.9022 & 0.9036 & 146062 & 7.521 \\
 & XGBoost\_Cmatched      & 0.9753 & 0.8926 & 0.8951 & 0.8951 & 939  & 8.273 \\
 & kNN\_5                 & 0.9733 & 0.8799 & 0.8752 & 0.8848 & 190624 & 15527.080 \\
 & GLADE+FPTM             & 0.9707 & 0.8727 & 0.8787 & 0.8828 & 41.5 & 5.471 \\
 & RandomForest\_Cmatched & 0.9611 & 0.8390 & 0.8631 & 0.8367 & 525  & 5.914 \\
 & LogisticRegression     & 0.9413 & 0.7763 & 0.7891 & 0.7853 & 4.8  & 0.901 \\
 & LinearSVM              & 0.8538 & 0.5146 & 0.6990 & 0.4660 & 4.7  & 0.823 \\
 & NaiveBayes             & 0.7673 & 0.4218 & 0.4803 & 0.5192 & 6.8  & 2.751 \\
\hline

\multirow{9}{*}{WUSTL-EHMS-2020}
 & GLADE+FPTM             & 0.9442 & 0.7976 & 0.9050 & 0.7645 & 25.2 & 1.446 \\
 & kNN\_5                 & 0.9375 & 0.7922 & 0.8445 & 0.7684 & 2857 & 364.308 \\
 & XGBoost                & 0.9406 & 0.7590 & 0.9151 & 0.7314 & 615  & 5.537 \\
 & XGBoost\_Cmatched      & 0.9354 & 0.7118 & 0.9198 & 0.6994 & 388  & -- \\
 & RandomForest           & 0.9320 & 0.6778 & 0.8910 & 0.6807 & 14385 & 4.380 \\
 & RandomForest\_Cmatched & 0.9305 & 0.6554 & 0.9671 & 0.6692 & 381  & -- \\
 & LinearSVM              & 0.9286 & 0.6449 & 0.6253 & 0.6657 & 2.4  & 0.122 \\
 & LogisticRegression     & 0.9280 & 0.6430 & 0.6220 & 0.6655 & 2.5  & 0.128 \\
 & NaiveBayes             & 0.1293 & 0.3643 & 0.6570 & 0.6668 & 3.0  & 0.511 \\
\hline
\end{tabular}
\end{table*}

\subsection{Serialization and Compression}
\label{sec:res-compression}

For deployment on resource-constrained edge devices, the trained
GLADE+FPTM pipeline is serialized into a single self-describing
container that stores both the GLADE binarization thresholds and
the FPTM clause bitmasks. Two equivalent formats containing identical
information are evaluated:

\begin{itemize}
    \item \textbf{\texttt{model.pkl}}: an uncompressed Python pickle
    file that stores the complete GLADE state (feature indices and
    thresholds) together with the TM rule representation, including
    the per-class positive and negative clause masks and their
    corresponding clamp values.

    \item \textbf{\texttt{model.fbz}}: the proposed
    \emph{Fuzzy-Bitmask Zstd (FBZ)} container, a compact binary format
    introduced in this work for deployment-oriented storage of the
    GLADE+FPTM inference artifact. The FBZ file consists of a fixed
    12-byte header followed by a Zstandard-compressed payload
    (zstd level~1). The header encodes the magic identifier, format
    version, compression level, dimensionality parameters
    (\texttt{n\_bits} and \texttt{n\_classes}), and the global clamp
    cap. Decoding requires only a standard zstd decompressor and
    little-endian integer unpacking, enabling portability across
    lightweight runtimes implemented in C, Rust, Go, or JavaScript.
\end{itemize}

Fig.~\ref{fig:storage_bar} compares the on-disk storage requirements
of the two formats. Across all four datasets, the FBZ container
reduces the deployable model footprint to approximately
21--30\% of the uncompressed pickle size. Specifically, the model
size is reduced from $25.2$\,KB to $7.1$\,KB on WUSTL-EHMS-2020,
from $51.1$\,KB to $13.7$\,KB on NSL-KDD, from $76.2$\,KB to
$16.2$\,KB on TON\_IoT, and from $41.5$\,KB to $12.3$\,KB on
MedSec-25. The compression process is strictly lossless, and exact
prediction equivalence after serialization and deserialization was
verified on the test split of every dataset.

\begin{figure}[t]
\centering
\includegraphics[width=\columnwidth]{TII-Articles-LaTeX-template/fig_storage_bar.pdf}
\caption{Storage footprint of the GLADE+FPTM model using
uncompressed (\texttt{model.pkl}) and compressed
(\texttt{model.fbz}) formats. Lower is better.}
\label{fig:storage_bar}
\end{figure}
\subsection{Compression Method Comparison}
\label{sec:res-codecs}

To position zstd-1 within the broader compression design space,
the same uncompressed FBZ payload was evaluated using six
lossless codecs spanning three compression families: zstandard
(zstd-1, zstd-5, and zstd-22), DEFLATE-based \textsc{gzip}-9,
LZMA2-based \texttt{lzma-9}, and \texttt{brotli-11}. To avoid
disk I/O overhead, all codecs operated on in-memory
\texttt{BytesIO} buffers such that the reported timings reflect
only codec execution cost. Each reported value corresponds to the
mean of 200 iterations, except for the computationally intensive
settings zstd-22, lzma-9, and brotli-11, which were averaged over
20 iterations. Compression effectiveness is evaluated using the
compression ratio (raw/compressed), while compression and
decompression latencies are reported in microseconds on a
logarithmic scale.

Three metrics are evaluated for each codec--dataset pair:
compression ratio (Fig.~\ref{fig:codec_size}), compression time
(Fig.~\ref{fig:codec_comp}), and decompression time
(Fig.~\ref{fig:codec_decomp}). The results indicate the following:

\begin{enumerate}
    \item \textbf{Compression ratio:} High-effort codecs
    (zstd-22, lzma-9, and brotli-11) achieve comparable compression
    performance across all datasets. zstd-1 provides slightly lower
    compression efficiency while remaining close to the
    best-performing methods.

    \item \textbf{Compression time:} zstd-1 consistently achieves
    the lowest compression latency, requiring only tens of
    microseconds per model. In contrast, lzma-9 and brotli-11 incur
    substantially higher encoding overhead.

    \item \textbf{Decompression time:} The zstd family provides the
    fastest decompression performance across all datasets. This
    characteristic is important for edge deployment scenarios,
    where the model is typically decompressed once during system
    initialization and reused for repeated inference.
\end{enumerate}

Based on these results, zstd-1 is selected as the default codec for
the FBZ container. Although higher compression levels provide
slightly improved storage reduction, zstd-1 offers the best trade-off
between model size, compression latency, and decompression
efficiency for practical edge deployment.

\begin{figure}[t]
\centering
\includegraphics[width=\columnwidth]{TII-Articles-LaTeX-template/fig_codec_size.pdf}
\caption{Compression ratio achieved by different lossless codecs
for the serialized GLADE+FPTM model. Higher values indicate better
compression efficiency.}
\label{fig:codec_size}
\end{figure}

\begin{figure}[t]
\centering
\includegraphics[width=\columnwidth]{TII-Articles-LaTeX-template/fig_codec_comp.pdf}
\caption{Compression latency of the evaluated codecs measured using
in-memory \texttt{BytesIO} buffers. Lower values indicate faster
compression.}
\label{fig:codec_comp}
\end{figure}

\begin{figure}[t]
\centering
\includegraphics[width=\columnwidth]{TII-Articles-LaTeX-template/fig_codec_size.pdf}
\caption{Decompression latency of the evaluated codecs under the
same experimental setup as Fig.~\ref{fig:codec_comp}. Lower values
indicate faster decompression.}
\label{fig:codec_decomp}
\end{figure}

\subsection{Edge Deployment: Binarisation Latency on Raspberry Pi 5}
\label{sec:res-edge}

To assess real-time feasibility on resource-constrained hardware, the
GLADE binarisation step was deployed on a Raspberry~Pi~5 and timed
against the TMU \texttt{StandardBinarizer} and a conventional
thermometer encoder. Two execution modes are reported: an element-wise
transform that emits one Boolean per literal, and a bit-packed
transform that produces 64-bit words ready for the FPTM clause
kernels. All values are mean per-sample latencies (in microseconds)
measured on the Raspberry~Pi~5 using its NEON SIMD path; speedups are
relative to GLADE. Results are summarised in
Table~\ref{tab:edge_latency}.

\begin{table}[t]
\centering
\caption{Per-sample binarisation latency on a Raspberry~Pi~5
(NEON SIMD path). Lower latency is better; speedups are computed
relative to GLADE.}
\label{tab:edge_latency}
\renewcommand{\arraystretch}{1.2}
\small
\begin{tabular}{|l|c|c|c|c|c|}
\hline
\textbf{Dataset} & \textbf{GLADE} & \textbf{Std.} & \textbf{Speedup} & \textbf{Thermo.} & \textbf{Speedup} \\
                 & ($\mu$s)       & ($\mu$s)      & vs Std.          & ($\mu$s)         & vs Thermo. \\
\hline
\multicolumn{6}{|l|}{\textit{Element-wise transform}} \\
\hline
WUSTL-EHMS-2020 & 0.162 & 0.905 & 5.6$\times$  & 0.459 & 2.8$\times$  \\
\hline
NSL-KDD         & 0.220 & 3.829 & 17.4$\times$ & 1.803 & 8.2$\times$  \\
\hline
TON\_IoT        & 0.094 & 1.604 & 17.1$\times$ & 0.577 & 6.1$\times$  \\
\hline
MedSec-25       & 0.287 & 3.477 & 12.1$\times$ & 0.906 & 3.2$\times$  \\
\hline
\multicolumn{6}{|l|}{\textit{Bit-packed transform}} \\
\hline
WUSTL-EHMS-2020 & 0.163 & 1.096 & 6.7$\times$  & 0.500 & 3.1$\times$  \\
\hline
NSL-KDD         & 0.219 & 4.200 & 19.2$\times$ & 2.373 & 10.8$\times$ \\
\hline
TON\_IoT        & 0.093 & 1.772 & 19.1$\times$ & 0.738 & 7.9$\times$  \\
\hline
MedSec-25       & 0.229 & 4.011 & 17.5$\times$ & 1.102 & 4.8$\times$  \\
\hline
\end{tabular}
\end{table}

On the Raspberry~Pi~5, GLADE binarisation completes in well under one
microsecond per sample on every dataset. Relative to the
\texttt{StandardBinarizer} it is $5.6\times$--$19.2\times$ faster, with
the largest gains on the higher-dimensional NSL-KDD and TON\_IoT
feature spaces, and it is up to $10.8\times$ faster than thermometer
encoding. The advantage is preserved in the bit-packed mode, confirming
that the compact Boolean representation produced by GLADE translates
directly into lower inference latency on embedded hardware.


% FIX 10: Added skeleton Conclusion section so \ref{sec:conclusion} resolves
\section{Conclusion}
\label{sec:conclusion}

This paper presented GLADE, a gap-aware lightweight adaptive
discretisation engine for transforming heterogeneous IoT feature
streams into compact Boolean representations, and its integration
with the Fuzzy Pattern Tsetlin Machine (FPTM) for resource-constrained
intrusion detection. Experimental results on four benchmark datasets
(NSL-KDD, TON\_IoT, MedSec-25, and WUSTL-EHMS-2020) demonstrated
that GLADE reduces the Boolean feature space by 45\%--65\% compared
to thermometric encoding while maintaining competitive or superior
classification performance. The combined GLADE+FPTM framework
achieves serialized model sizes as low as 7.1\,KB and, on a
Raspberry~Pi~5 edge platform, performs binarisation up to 19$\times$
faster than the standard binariser and up to 11$\times$ faster than
thermometer encoding, confirming its suitability for real-time
deployment in memory- and compute-constrained IoT environments.

\bibliographystyle{IEEEtran}
\bibliography{references}

\end{document}